\chapter{12.11-1 等腰三角形性质}


\section*{二、填空题（共6小题）}


\blanklist[7]

\item 如图，阴影部分是边长为1的小正三角形，$A, B, C, D, E, F, G, H$ 分别是8个正三角形，则 $A$ 和 $B$ 的边长分别是（\quad ）
  \par\smallskip\noindent
  \begin{minipage}[t]{0.44\linewidth}
  \vspace{0pt}
 \optionlist
  \item 2, 4
  \item 2.5, 5
  \item 3, 6
  \item 4, 8
  \end{enumerate}
  \end{minipage}
  \hfill
  \begin{minipage}[t]{0.2\linewidth}
  \questionimage{12.11-7.png}
  \end{minipage}


\item 图①是一块边长为1，周长记为 $P_1$ 的正三角形纸板，沿图①的底边剪去一块边长为 $\frac{1}{2}$ 的正三角形纸板后得到图②，然后沿同一底边依次剪去一块更小的正三角形纸板（即其边长为前一块被剪掉正三角形纸板边长的 $\frac{1}{2}$）后，得图③，④，…，记第 $n$（$n \geq 3$）块纸板的周长为 $P_n$，则 $P_n - P_{n-1}$ 的值为（    ）
  \par\smallskip\noindent
  \begin{minipage}[t]{0.44\linewidth}
  \vspace{0pt}
  \optionlist
  \item $\left( \frac{1}{4} \right)^{n-1}$
  \item $\left( \frac{1}{4} \right)^n$
  \item $\left( \frac{1}{2} \right)^{n-1}$
  \item $\left( \frac{1}{2} \right)^n$
  \end{enumerate}
    \end{minipage}
  \hfill
  \begin{minipage}[t]{0.55\linewidth}
  \questionimage{12.11-8.png}
  \end{minipage}

% 插入图片：此处应有图①、②、③、④

\end{enumerate}

\section*{二、填空题（共6小题）}

\blanklist[9]

\item 如图，已知 $\triangle ABC$ 是等边三角形，点 $O$ 是 $BC$ 上任意一点，$OE, OF$ 分别与两边垂直，等边三角形的高为5，则 $OE + OF$ 的值为\_\_\_\_\_\_。

\rightquestionimage[width=0.35\textwidth]{0.35\textwidth}{12.11-9.png}

\item 三个等边三角形的位置如图所示，若 $\angle 3 = 40^\circ$，则 $\angle 1 + \angle 2 =$\_\_\_\_\_\_。

\rightquestionimage[width=0.35\textwidth]{0.35\textwidth}{12.11-10.png}

\item 如图，$\triangle ABC$ 是边长为3的等边三角形，$\triangle BDC$ 是等腰三角形，且 $\angle BDC = 120^\circ$。以 $D$ 为顶点作一个 $60^\circ$ 角，使其两边分别交 $AB$ 于点 $M$，交 $AC$ 于点 $N$，连接 $MN$，则 $\triangle AMN$ 的周长为\_\_\_\_\_\_。

\rightquestionimage[width=0.35\textwidth]{0.35\textwidth}{12.11-11.png}

\item 将一个平面图形分成面积相等的两部分的直线叫做该平面图形的“面线”，“面线”被这个平面图形截得的线段叫做该图形的“面径”，例如圆的直径就是它的“面径”。已知等边三角形的边长为2，则它的“面径”长可以是\_\_\_\_\_\_（写出2个）。


\item 如图，将边长为1的正三角形 $OAP$ 沿 $x$ 轴正方向连续无滑动的翻转2019次，点 $P$ 依次落在点 $P_1, P_2, \ldots, P_{2019}$ 的位置，则点 $P_{2019}$ 的横坐标为\_\_\_\_\_\_。

\rightquestionimage[width=0.35\textwidth]{0.35\textwidth}{12.11-13.png}

\item 如图，一只蜘蛛结了一张很大的网，直线 $MN, PQ$ 与 $x$ 轴所夹的锐角都为 $60^\circ$，第1个结点 $A_1$ 在原点，此后各个结点均按逆时针排列，同一直线上相邻两个结点之间的距离都是1个单位长，那么第91个结点 $A_{91}$ 的坐标为\_\_\_\_\_\_。

\rightquestionimage[width=0.35\textwidth]{0.35\textwidth}{12.11-14.png}

\end{enumerate}

\newpage

\section*{三、解答题（共3小题）}

\solutionlist[15]

\item 如图，以 $\triangle ABC$ 的两边 $AB, AC$ 向外作等边三角形 $ABE$ 和等边三角形 $ACD$，连接 $BD, CE$，交于 $O$。

(1) 试写出图中和 $BD$ 相等的一条线段并说明你的理由；

(2) 求出 $BD$ 和 $CE$ 的夹角大小，若改变 $\triangle ABC$ 的形状，这个夹角的度数会发生变化吗？请说明理由。

\rightquestionimage[width=0.35\textwidth]{0.35\textwidth}{12.11-15.png}

\vspace{3cm}

\item 将两个等边 $\triangle ABC$ 和 $\triangle DEF$（$DE > AB$）如图所示摆放，点 $D$ 是 $BC$ 上的一点（除 $B, C$ 点外）。把 $\triangle DEF$ 绕顶点 $D$ 顺时针旋转一定的角度，使得边 $DE, DF$ 与 $\triangle ABC$ 的边（除 $BC$ 边外）分别相交于点 $M, N$。

(1) $\angle BMD$ 和 $\angle CDN$ 相等吗？

(2) 画出使 $\angle BMD$ 和 $\angle CDN$ 相等的所有情况的图形；

(3) 在 (2) 题中任选一种图形说明 $\angle BMD$ 和 $\angle CDN$ 相等的理由。


\rightquestionimage[width=0.35\textwidth]{0.35\textwidth}{12.11-16.png}

\end{enumerate}

\newpage

\solutionlist[17]

\item 已知：等边三角形 $ABC$，

(1) 如图1，$P$ 为等边 $\triangle ABC$ 外一点，且 $\angle BPC = 120^\circ$。试猜想线段 $BP$、$PC$、$AP$ 之间的数量关系，并证明你的猜想；

(2) 如图2，$P$ 为等边 $\triangle ABC$ 内一点，且 $\angle APD = 120^\circ$。求证：$PA + PD + PC > BD$。

\centerquestionimage[width=0.70\textwidth]{0.70\textwidth}{12.11-17.png}


\end{enumerate}
